\documentclass[12pt,a4paper]{book}
\usepackage[utf8]{inputenc}
\usepackage[T1]{fontenc}
\usepackage[spanish]{babel}
\usepackage{amsmath, amssymb, amsfonts}
\usepackage{graphicx}
\usepackage{geometry}
\usepackage{booktabs}
\usepackage{enumitem}
\usepackage{tcolorbox}

\geometry{margin=2.5cm}

% Entornos destacados instruccionales
\newtcolorbox{definicion}[1][]{
    colback=blue!5!white,
    colframe=blue!75!black,
    fonttitle=\bfseries,
    title=Definición #1
}

\newtcolorbox{ejemplo}[1][]{
    colback=green!5!white,
    colframe=green!50!black,
    fonttitle=\bfseries,
    title=Ejemplo #1
}

\begin{document}

% ----------------------------------------------------------------------
% FRONTMATTER
% ----------------------------------------------------------------------
\frontmatter

\begin{titlepage}
\begin{center}
    \vspace*{2cm}
    \textbf{\Large UNIVERSIDAD AUTÓNOMA GABRIEL RENÉ MORENO} \\[0.5cm]
    \textbf{\large FACULTAD DE CIENCIAS CONTABLES} \\[0.5cm]
    \textbf{\large CARRERA DE CONTADURÍA PÚBLICA} \\[2cm]
    
    \textbf{\Huge ESTADÍSTICA II MAT-260} \\[1.5cm]
    \textbf{\Large Texto Guía de Probabilidades e Inferencia} \\[2cm]
    
    \textbf{\large Docente: MSc. Anselmo Salguero Arano} \\[2cm]
    \vfill
    \textbf{\large Santa Cruz - Bolivia} \\
\end{center}
\end{titlepage}

\tableofcontents

% ----------------------------------------------------------------------
% MAINMATTER
% ----------------------------------------------------------------------
\mainmatter

% ----------------------------------------------------------------------
% UNIDAD 1
% ----------------------------------------------------------------------
\chapter{Unidad N°1: Análisis Combinatorio}

\section{Introducción}
El análisis combinatorio trata de los diversos tipos de arreglos o selecciones que se pueden formar con los elementos de un conjunto[cite: 2]. Tomando en cuenta las propiedades de los grupos distintos que se pueden formar, diferenciándose entre sí por orden de colocación, por orden según la clase de elementos, y por el número de elementos que entra en cada clase[cite: 2]. El estudio del análisis combinatorio nos servirá como base para poder resolver y comprender problemas sobre probabilidades[cite: 2].

\section{Conteo}
El proceso de conteo simplifica la cuantificación de eventos[cite: 2]. En el estudio de lo que es posible existen esencialmente dos tipos de problemas: citar todo lo que puede suceder en una situación dada y determinar cuántas cosas diferentes pueden acontecer[cite: 2]. 

\subsection{Leyes Fundamentales}
\begin{definicion}[Ley de la Multiplicación]
Si una elección consta de $k$ pasos, donde el primero puede realizarse de $n_{1}$ maneras, el segundo de $n_{2}$ maneras, y así sucesivamente, el número total de formas en que toda la situación puede resolverse es el producto de todas estas formas[cite: 2].
\[
\text{Total de formas} = n_{1} \times n_{2} \times \dots \times n_{k}
\]
Aplica para eventos que no son mutuamente excluyentes[cite: 2].
\end{definicion}

\begin{definicion}[Ley de la Adición]
Si un evento $A$ puede ocurrir de $n$ formas diferentes y un evento $B$ de $m$ formas diferentes, y ambos son eventos mutuamente excluyentes (la ocurrencia de uno impide la del otro), el número total de formas es la suma de sus formas individuales[cite: 2].
\[
\text{Total de formas} = n + m
\]
\end{definicion}

\section{Permutaciones}
El número de permutaciones de $n$ objetos es el número de formas en las que pueden acomodarse esos objetos en términos de orden[cite: 2].

\textbf{Permutaciones de $n$ objetos (Factorial):} El número total de arreglos posibles de $n$ objetos es $n$ factorial[cite: 2], denotado por $P_{n} = n!$.
\[
P_{n} = n! = n \times (n-1) \times (n-2) \times \dots \times 3 \times 2 \times 1
\]
\textbf{Permutaciones de $n$ objetos tomados de $r$ a la vez:}
\[
_{n}P_{r} = \frac{n!}{(n-r)!}
\]

\section{Combinaciones}
En el caso de las combinaciones, lo importante es el número de agrupaciones diferentes de objetos que pueden ocurrir sin importar el orden[cite: 2].
\[
_{n}C_{r} = \frac{n!}{r! (n-r)!}
\]

\begin{ejemplo}[Diferencia entre Permutación y Combinación]
A cada grupo ordenado de elementos se le llama permutación, mientras que a un grupo no ordenado de elementos se le llama combinación[cite: 2].
El conjunto $\{A, B, C\}$ tiene 6 permutaciones: $(A,B,C)$, $(A,C,B)$, $(B,A,C)$, $(B,C,A)$, $(C,A,B)$, $(C,B,A)$, pero representa 1 sola combinación[cite: 2].
\end{ejemplo}

\section{Ejercicios propuestos}
\begin{enumerate}
    \item Determine el número de permutaciones posibles con las letras de la palabra "UAGRM".
    \item Un comité de 4 personas debe formarse a partir de un grupo de 10 contadores. ¿De cuántas formas se puede seleccionar el comité?
    \item Un estudiante debe responder 8 de 10 preguntas en un examen. ¿De cuántas maneras puede elegirlas?
\end{enumerate}

% ----------------------------------------------------------------------
% UNIDAD 2
% ----------------------------------------------------------------------
\chapter{Unidad N°2: Probabilidad}

\section{Introducción a la Probabilidad}
Experimento es el proceso por el cual se obtiene un resultado de una observación[cite: 2]. Puede ser determinístico (el resultado se puede predecir en forma precisa) o aleatorio (el resultado no se puede predecir con exactitud)[cite: 2].

\begin{definicion}[Espacio Muestral y Evento]
Un conjunto $S$ que consta de todos los resultados posibles de un experimento aleatorio se llama espacio muestral, y cada resultado se llama punto muestral[cite: 2]. Un evento ($N$) es un subconjunto cualquiera de un espacio muestral[cite: 2].
\end{definicion}

\section{Definición de Probabilidad}
Históricamente se han desarrollado tres enfoques conceptuales para definir probabilidad y para determinar valores de probabilidad[cite: 2]:
\begin{itemize}
    \item \textbf{Definición Clásica (A priori):} Se basa en la suposición de que cada uno de los resultados es igualmente probable[cite: 2]. $P(A) = \frac{N(A)}{S}$[cite: 2].
    \item \textbf{Frecuencia Relativa (Empírica):} Se determina con base en la proporción de veces que ocurre un resultado favorable en un determinado número de observaciones[cite: 2].
    \item \textbf{Subjetiva:} Es el grado de confianza que una persona tiene en que el evento ocurra, con base en toda la evidencia que tiene disponible[cite: 2].
\end{itemize}

\section{Reglas de Adición y Multiplicación}
La suma de la probabilidad de ocurrencia más la probabilidad de no ocurrencia siempre es igual a 1[cite: 2]:
\[
P(A) + P(A') = 1
\]
\textbf{Regla general de adición para eventos no excluyentes:}
\[
P(A \cup B) = P(A) + P(B) - P(A \cap B)
\]

\begin{definicion}[Probabilidad Condicional]
Cuando dos eventos son dependientes, se utiliza el concepto de probabilidad condicional para designar la ocurrencia de un evento relacionado[cite: 2].
\[
P(A|B) = \frac{P(A \cap B)}{P(B)}
\]
\end{definicion}

\section{Teorema de Bayes}
Se aplica en el contexto de eventos secuenciales[cite: 2]. La fórmula de cálculo para el teorema de Bayes es[cite: 2]:
\[
P(A|B) = \frac{P(A) \cdot P(B|A)}{P(A) \cdot P(B|A) + P(A') \cdot P(B|A')}
\]

\section{Ejercicios propuestos}
\begin{enumerate}
    \item En la Facultad de Ciencias Contables, la probabilidad de que un estudiante apruebe Estadística II es 0.70. ¿Cuál es la probabilidad de que no apruebe?
    \item Si se lanzan dos dados, ¿cuál es la probabilidad de que la suma de sus caras sea 7?
    \item Aplicar el teorema de Bayes para una fábrica con tres máquinas que producen el 20\%, 30\% y 50\% de las calculadoras, con tasas de defecto del 5\%, 4\% y 2\% respectivamente.
\end{enumerate}


% ----------------------------------------------------------------------
% UNIDAD 3
% ----------------------------------------------------------------------
\chapter{Unidad N°3: Distribuciones de Probabilidad para Variables Aleatorias Discretas}

\section{Distribución Binomial}
A un número fijo $n$ de repeticiones independientes de una prueba de Bernoulli se llama experimento binomial[cite: 2]. Solo son dos posibles resultados mutuamente excluyentes en cada ensayo u observación: éxito ($p$) o fracaso ($q$)[cite: 2].
\[
P(x|n,p) = \frac{n!}{x!(n-x)!} p^x q^{n-x}
\]
Donde $x$ es el número específico de éxitos y $n$ es el número de observaciones o ensayos[cite: 2].

\begin{ejemplo}[Medidas de tendencia central y dispersión]
Para la distribución binomial, el valor esperado o media es $E(x) = n \times p$ y la varianza es $V(x) = n \times p \times q = \sigma^{2}$[cite: 2].
\end{ejemplo}

\section{Distribución Hipergeométrica}
Cuando el muestreo se realiza sin reemplazo para cada uno de los elementos que se toman de una población finita, existe un cambio sistemático en la probabilidad de éxitos[cite: 2].
\[
P(X|N,T,n) = \frac{ \binom{T}{x} \binom{N-T}{n-x} }{ \binom{N}{n} }
\]
Donde $N$ es el número de elementos de la población y $T$ es el número total de éxitos incluidos en la población[cite: 2].

\section{Distribución de Poisson}
Puede utilizarse la distribución de Poisson para determinar la probabilidad de que ocurra un número asignado de eventos cuando estos ocurren en un continuo de tiempo y espacio[cite: 2].
\[
P(x|\lambda) = \frac{e^{-\lambda} \lambda^x}{x!}
\]
El valor esperado y la varianza son iguales a la media de la distribución: $E(x) = V(x) = \lambda$[cite: 2].

\section{Ejercicios propuestos}
\begin{enumerate}
    \item Calcule la probabilidad de obtener exactamente 3 caras al lanzar una moneda 5 veces.
    \item Una caja contiene 10 artículos, 3 defectuosos. Si se extraen 4 sin reemplazo, determine la probabilidad de obtener 1 defectuoso.
    \item El número promedio de clientes que llega a un banco en un minuto es 2. ¿Cuál es la probabilidad de que lleguen exactamente 3 en un minuto?
\end{enumerate}


% ----------------------------------------------------------------------
% UNIDAD 4
% ----------------------------------------------------------------------
\chapter{Unidad N°4: Distribución de Probabilidad para Variables Aleatorias Continuas}

\section{Introducción}
Existen diversas distribuciones continuas de probabilidad comunes que son aplicables como modelo a una amplia gama de variables continuas[cite: 2]. Una densidad de probabilidad se caracteriza por el hecho de que el área situada debajo de la curva entre dos valores $a$ y $b$, da la probabilidad de que una variable aleatoria tome un valor contenido en el intervalo $(a,b)$[cite: 2]. El área total situada debajo de la curva siempre es igual a 1[cite: 2].

\section{Distribución Normal o Estándar}
La gráfica de una distribución normal es una curva en forma de campana que se extiende indefinidamente en ambos sentidos[cite: 2]. 

\begin{definicion}[Estandarización]
Se utiliza para transformar cualquier distribución normal en la distribución normal estándar[cite: 2]. La fórmula es:
\[
Z = \frac{x - \mu}{\sigma}
\]
Donde $\mu$ es la media y $\sigma$ es la desviación estándar[cite: 2].
\end{definicion}

Las distribuciones de estadística como la media muestral y la proporción muestral tienen distribución normal cuando el tamaño de la muestra es grande, sin importar la forma de la distribución de la población de origen[cite: 2].

\section{Ejercicios propuestos}
\begin{enumerate}
    \item Sabiendo que $Z$ sigue una distribución normal estándar, calcule el área bajo la curva a la derecha de $Z = 1.25$.
    \item Los sueldos de los contadores tienen distribución normal con media 5000 Bs y desviación estándar 800 Bs. ¿Qué porcentaje gana menos de 4000 Bs?
    \item Encuentre el valor de $Z$ que deja un área de 0.95 a su izquierda.
\end{enumerate}


% ----------------------------------------------------------------------
% UNIDAD 5
% ----------------------------------------------------------------------
\chapter{Unidad N°5: Distribuciones de Muestreo e Intervalos de Confianza para la Media}

\section{Población y Muestra}
Llevar a cabo algunos estudios sobre toda una población resulta casi imposible o impráctico, por lo que la solución es llevar a cabo el estudio basándose en un subconjunto de ésta denominada muestra[cite: 2].
\begin{itemize}
    \item \textbf{Población:} es el conjunto de todos los elementos que interesan en un estudio[cite: 2].
    \item \textbf{Muestra:} es un subconjunto de la población[cite: 2].
\end{itemize}

\section{Estimación Puntual y Teorema del Límite Central}
Un estimador puntual es la estimación de un parámetro poblacional, basado en un solo número[cite: 2]. 

\begin{definicion}[Teorema del Límite Central]
Al aumentar el tamaño de la muestra, la distribución muestral de la media se aproxima a la forma de la distribución normal sin importar la forma de la distribución de las mediciones individuales de la población[cite: 2]. Puede suponerse que la distribución muestral de la media es aproximadamente normal cuando el tamaño de la muestra es $n \ge 30$[cite: 2].
\end{definicion}

\section{Intervalos de Confianza para la Media}
Cuando puede utilizarse la distribución normal de probabilidad, el intervalo de confianza para la media se determina mediante[cite: 2]:
\[
\overline{X} \pm Z \sigma_{\overline{x}}
\]

\textbf{Determinación del tamaño de la muestra:} Si se conoce $\sigma$ o puede estimarse, el tamaño de la muestra que se requiere con base en la distribución normal es[cite: 2]:
\[
n = \left[ \frac{Z \sigma}{E} \right]^2
\]
Donde $E$ es el factor de error máximo permitido en el intervalo[cite: 2].

\section{Ejercicios propuestos}
\begin{enumerate}
    \item Calcule el intervalo de confianza del 95\% para una muestra de 36 elementos con media 120 y desviación estándar 15.
    \item Se requiere estimar la media poblacional con un margen de error de 5 y nivel de confianza del 99\%. Si $\sigma = 20$, ¿cuál debe ser el tamaño de la muestra?
    \item Explique la importancia del Teorema del Límite Central.
\end{enumerate}

% ----------------------------------------------------------------------
% BACKMATTER
% ----------------------------------------------------------------------
\backmatter

\chapter{Referencias}
\begin{itemize}
    \item Salguero Arano, A. (2026). \textit{TEXTO ESTADÍSTICA II: Estadística inferencial}. Santa Cruz - Bolivia: Universidad Autónoma Gabriel René Moreno[cite: 2].
    \item López, A., y Burgoa, C. (2020). \textit{Estadística Aplicada a los Negocios}.
\end{itemize}

\end{document}